\documentclass[11pt]{article}

% Change "review" to "final" to generate the final (sometimes called camera-ready) version.
% Change to "preprint" to generate a non-anonymous version with page numbers.
\usepackage[]{acl}
\usepackage{times}
\usepackage{latexsym}

% For proper rendering and hyphenation of words containing Latin characters (including in bib files)
\usepackage[T1]{fontenc}
% For Vietnamese characters
% \usepackage[T5]{fontenc}
% See https://www.latex-project.org/help/documentation/encguide.pdf for other character sets

% This assumes your files are encoded as UTF8
\usepackage[utf8]{inputenc}
\usepackage{microtype}
\usepackage{inconsolata}
\usepackage{graphicx}

% If the title and author information does not fit in the area allocated, uncomment the following
%
%\setlength\titlebox{<dim>}
%
% and set <dim> to something 5cm or larger.

\title{Group 21 Progress Report:\\My Group's Project Name}


\author{First Author, Second Author, Third Author \\
  \texttt{\{parasutt,macid2,macid3\}@mcmaster.ca} }

%\author{
%  \textbf{First Author\textsuperscript{1}},
%  \textbf{Second Author\textsuperscript{1,2}},
%  \textbf{Third T. Author\textsuperscript{1}},
%  \textbf{Fourth Author\textsuperscript{1}},
%\\
%  \textbf{Fifth Author\textsuperscript{1,2}},
%  \textbf{Sixth Author\textsuperscript{1}},
%  \textbf{Seventh Author\textsuperscript{1}},
%  \textbf{Eighth Author \textsuperscript{1,2,3,4}},
%\\
%  \textbf{Ninth Author\textsuperscript{1}},
%  \textbf{Tenth Author\textsuperscript{1}},
%  \textbf{Eleventh E. Author\textsuperscript{1,2,3,4,5}},
%  \textbf{Twelfth Author\textsuperscript{1}},
%\\
%  \textbf{Thirteenth Author\textsuperscript{3}},
%  \textbf{Fourteenth F. Author\textsuperscript{2,4}},
%  \textbf{Fifteenth Author\textsuperscript{1}},
%  \textbf{Sixteenth Author\textsuperscript{1}},
%\\
%  \textbf{Seventeenth S. Author\textsuperscript{4,5}},
%  \textbf{Eighteenth Author\textsuperscript{3,4}},
%  \textbf{Nineteenth N. Author\textsuperscript{2,5}},
%  \textbf{Twentieth Author\textsuperscript{1}}
%\\
%\\
%  \textsuperscript{1}Affiliation 1,
%  \textsuperscript{2}Affiliation 2,
%  \textsuperscript{3}Affiliation 3,
%  \textsuperscript{4}Affiliation 4,
%  \textsuperscript{5}Affiliation 5
%\\
%  \small{
%    \textbf{Correspondence:} \href{mailto:email@domain}{email@domain}
%  }
%}

\begin{document}
\maketitle
% \begin{abstract}
% \end{abstract}

\section{Introduction}

% Here, write a brief introduction to the problem you are solving. This can be adapted from your problem description and motivation from the original proposal. This should be around 0.25-0.5 pages.
Identifying which instrument produced a sound is harder than it seems. Instruments in the same family, such as the oboe and cor anglais or the violin and viola, have very similar timbres, and a single instrument can sound quite different depending on the pitch, dynamics, and articulation it is played with. Being able to identify an instrument from its sound alone would make recordings easier to search and catalogue, for example to find other recordings that feature the same instrument as a given composition.

In this project, we build a machine learning system that classifies orchestral instruments from short audio recordings. We frame this as a single-label, multi-class classification task over 19 brass, woodwind, and string instruments, using single-note clips from the Philharmonia Orchestra sound sample library \citep{philharmonia_samples}. Each clip is converted from a raw waveform into a time--frequency representation using the Fourier transform. From this, we either extract hand-crafted spectral features for classical models such as SVMs and random forests, or compute log-mel spectrograms as input to a small convolutional neural network.

The task has several challenges. Audio must first be processed with signal processing techniques before it can be used as model input. The model must learn timbre rather than shortcuts such as pitch or loudness, since some instruments (e.g., tuba, contrabassoon, and double bass) can be identified from pitch range alone. The dataset is also heavily imbalanced, with over 1,500 violin clips but fewer than 100 for some instruments. Finally, because the recordings are clean studio samples, a model trained on them may not generalize to real-world recordings.

\section{Related Work}

% Here, talk about the related work you encountered for your approach. Cite at least 5 references. Refer to item 2. No one has done exactly your task? Write about the most similar thing you can find. This should be around 0.25-0.5 pages.
Instrument recognition is a well-studied problem in music information retrieval. Early work focused on isolated notes, as in our task, and relied on hand-crafted features. \citet{eronen2000musical} combined cepstral coefficients with temporal features describing a note's attack and decay, and \citet{herrera2003automatic} surveyed the spectral and temporal features and classifiers (such as $k$-nearest neighbours, decision trees, and SVMs) used for classifying isolated instrument sounds. These works motivate our classical baseline of MFCC statistics, spectral descriptors, and attack time fed to an SVM or random forest.

More recent work has used larger datasets and deep learning. The NSynth dataset \citep{engel2017nsynth} contains over 300,000 four-second notes from roughly 1,000 instruments labelled by instrument family, although its main focus is audio synthesis rather than classification. The IRMAS dataset \citep{bosch2012irmas} targets the harder problem of recognizing the predominant instrument in polyphonic music across 11 classes. On IRMAS, \citet{han2017deep} showed that a convolutional neural network trained on mel spectrograms outperforms feature-based approaches, which motivates our log-mel spectrogram CNN. Finally, PANNs \citep{kong2020panns} are CNNs pretrained on large-scale general audio data that transfer well to other audio classification tasks, and could serve as a pretrained alternative if our small CNN is limited by the size of our dataset.

Our task is different from these, we are classifying orchestral instruments rather than their broad families, using clean single-note recordings. Making it more fine-grained than Nsynth's but simpler than polyphonic recognition. Because every instrument is recorded across many pitches and dynamics. 
\section{Dataset}


% You should write about your dataset here, following the guidelines regarding item 1. This section may be 0.5-1 pages. Depending on your specific dataset, you may want to include subsections for the preprocessing, annotation, etc.

We use the Philharmonia Orchestra sound sample library \citep{philharmonia_samples}, a freely available collection of recordings made by members of the Philharmonia Orchestra. The library is distributed as MP3 files grouped by instrument, and each file name records the instrument, pitch, duration, dynamic, and articulation of the clip (e.g., \texttt{flute\_A4\_025\_forte\_normal.mp3}). The instrument name serves as our class label, and the remaining fields let us analyse which factors affect our models, such as whether accuracy drops at extreme pitches or quiet dynamics, without any manual annotation.

We use 19 instruments from three families: four brass, eight woodwind, and seven string instruments. The strings include the plucked guitar, mandolin, and banjo as well as the bowed orchestral strings. Table~\ref{tab:dataset} lists the number of clips and the pitch range of each instrument. While inspecting the data, we found that the cor anglais folder contains every recording twice, once under \texttt{cor-anglais} and once under its alternate name \texttt{english-horn}, and that the two copies are byte-identical. We remove these duplicates, since otherwise copies of the same recording could appear in both the training and test sets. This leaves 13,535 clips.

The recordings vary along several axes. Most clips are single notes with a duration of 0.25, 0.5, 1, or 1.5 seconds (11,485 clips), with fewer long and very long sustained notes (961 clips) and 1,089 short phrases. Dynamics range from pianissimo to fortissimo. About 84\% of clips use the instrument's standard articulation, while the rest cover 49 other playing techniques such as pizzicato, staccato, trills, tremolo, and muted playing. Pitch ranges overlap heavily between instruments, but the contrabassoon and tuba reach down to A$\sharp$0 and the double bass to C1, below every other instrument in the dataset.

The class distribution is highly imbalanced. Violin has 1,502 clips, about twenty times as many as banjo (74), and the three plucked strings together make up less than 2\% of the data. As a result, overall accuracy alone can be misleading, since a model can score well while performing poorly on the rarest classes.

\begin{table}[t]
  \centering
  \small
  \begin{tabular}{lrl}
    \hline
    \textbf{Instrument} & \textbf{Clips} & \textbf{Pitch range} \\
    \hline
    \multicolumn{3}{l}{\textit{Brass}} \\
    French horn    & 652   & A$\sharp$1--F5 \\
    Trombone       & 831   & E2--F6 \\
    Trumpet        & 485   & E2--E6 \\
    Tuba           & 972   & A$\sharp$0--F4 \\
    \hline
    \multicolumn{3}{l}{\textit{Woodwind}} \\
    Bass clarinet  & 944   & C2--C6 \\
    Bassoon        & 720   & A$\sharp$1--G5 \\
    Clarinet       & 846   & D3--C7 \\
    Contrabassoon  & 710   & A$\sharp$0--D4 \\
    Cor anglais    & 691   & E3--B5 \\
    Flute          & 878   & C4--F7 \\
    Oboe           & 596   & A$\sharp$3--A$\sharp$6 \\
    Saxophone      & 733   & D3--G6 \\
    \hline
    \multicolumn{3}{l}{\textit{Strings}} \\
    Banjo          & 74    & C3--E6 \\
    Cello          & 889   & C2--C7 \\
    Double bass    & 852   & C1--G4 \\
    Guitar         & 106   & E2--E6 \\
    Mandolin       & 80    & G3--A6 \\
    Viola          & 974   & C3--F$\sharp$7 \\
    Violin         & 1,502 & G3--E8 \\
    \hline
    \textbf{Total} & \textbf{13,535} & \\
    \hline
  \end{tabular}
  \caption{\label{tab:dataset}
    Number of clips and pitch range for each instrument, after removing duplicate cor anglais recordings.
  }
\end{table}
\section{Features}


% Describe any features you used for your model, or how your data was input to your model. Are you doing feature engineering or feature selection? Are you learning embeddings? Is it all part of one neural network? Refer to item 2. This may range from 0.25 pages to 0.5 pages.
Both of our models share the same preprocessing. Each MP3 file is decoded, converted to mono, and resampled to 16\,kHz. We then trim leading and trailing silence and pad or crop every clip to a fixed length, which we will choose experimentally between 1 and 2 seconds. From this waveform, we compute a short-time Fourier transform, map its frequency bins onto the mel scale, which approximates how humans perceive pitch, and take the logarithm to compress the large range of energy between quiet and loud frequencies. The result is a log-mel spectrogram that shows how the energy at each frequency changes over the course of the note.

\paragraph{Classical baseline.} For the SVM and random forest, we engineer a fixed-length feature vector for each clip. Since these models cannot handle a variable-length sequence of frames, we summarize each frame-level feature by its mean and standard deviation over time. The features, following \citet{eronen2000musical} and \citet{herrera2003automatic}, are:
\begin{itemize}
  \item \textbf{MFCCs}, which describe the overall shape of the spectrum and are the main indicator of timbre;
  \item \textbf{spectral centroid and rolloff}, which measure how bright a sound is;
  \item \textbf{spectral flatness and zero-crossing rate}, which measure how noisy, rather than tonal, a sound is;
  \item \textbf{attack time}, the time the note takes to reach its peak amplitude, which separates plucked strings from bowed and blown instruments.
\end{itemize}
We standardize each feature using statistics from the training set. We do not perform explicit feature selection, since the feature set is small and both models are reasonably robust to uninformative features.

\paragraph{CNN.} The CNN takes the log-mel spectrogram directly as a single-channel image, with mel frequency on one axis and time on the other. It uses no hand-crafted features. Instead, the convolutional layers learn their own representation of the spectrogram, and the whole model, from spectrogram to class prediction, is trained end-to-end as one network. To reduce overfitting on our relatively small dataset, we plan to apply SpecAugment-style masking \citep{park2019specaugment}, which hides random bands of frequencies and time steps during training. If the small CNN proves too limited, we can replace its learned representation with embeddings from a pretrained PANNs model \citep{kong2020panns}.

\section{Implementation}

Describe your model and implementation here. Refer to item 4. This may take around a page.

\section{Results and Evaluation}

How are you evaluating your model? What results do you have so far? What are your baselines? Refer to item 5. This may take around 0.5 pages.

\section{Feedback and Plans}

Write about your plans for the remainder of the project. This should include a discussion of the feedback you received from your TA, and how you plan to improve your approach. Reflect on your implementation and areas for improvement. Refer to item 6. This may be around 0.5 pages.

\section{Template Notes}

You can remove this section or comment it out, as it only contains instructions for how to use this template. You may use subsections in your document as you find appropriate.

\subsection{Tables and figures}

See Table~\ref{citation-guide} for an example of a table and its caption.
See Figure~\ref{fig:experiments} for an example of a figure and its caption.


\begin{figure}[t]
  \includegraphics[width=\columnwidth]{example-image-golden}
  \caption{A figure with a caption that runs for more than one line.
    Example image is usually available through the \texttt{mwe} package
    without even mentioning it in the preamble.}
  \label{fig:experiments}
\end{figure}

\begin{figure*}[t]
  \includegraphics[width=0.48\linewidth]{example-image-a} \hfill
  \includegraphics[width=0.48\linewidth]{example-image-b}
  \caption {A minimal working example to demonstrate how to place
    two images side-by-side.}
\end{figure*}


\subsection{Citations}

\begin{table*}
  \centering
  \begin{tabular}{lll}
    \hline
    \textbf{Output}           & \textbf{natbib command} & \textbf{ACL only command} \\
    \hline
    \citep{Gusfield:97}       & \verb|\citep|           &                           \\
    \citealp{Gusfield:97}     & \verb|\citealp|         &                           \\
    \citet{Gusfield:97}       & \verb|\citet|           &                           \\
    \citeyearpar{Gusfield:97} & \verb|\citeyearpar|     &                           \\
    \citeposs{Gusfield:97}    &                         & \verb|\citeposs|          \\
    \hline
  \end{tabular}
  \caption{\label{citation-guide}
    Citation commands supported by the style file.
  }
\end{table*}

Table~\ref{citation-guide} shows the syntax supported by the style files.
We encourage you to use the natbib styles.
You can use the command \verb|\citet| (cite in text) to get ``author (year)'' citations, like this citation to a paper by \citet{Gusfield:97}.
You can use the command \verb|\citep| (cite in parentheses) to get ``(author, year)'' citations \citep{Gusfield:97}.
You can use the command \verb|\citealp| (alternative cite without parentheses) to get ``author, year'' citations, which is useful for using citations within parentheses (e.g. \citealp{Gusfield:97}).

\subsection{References}

\nocite{Ando2005,andrew2007scalable,rasooli-tetrault-2015}

Many websites where you can find academic papers also allow you to export a bib file for citation or bib formatted entry. Copy this into the \texttt{custom.bib} and you will be able to cite the paper in the \LaTeX{}. You can remove the example entries.

\subsection{Equations}

An example equation is shown below:
\begin{equation}
  \label{eq:example}
  A = \pi r^2
\end{equation}

Labels for equation numbers, sections, subsections, figures and tables
are all defined with the \verb|\label{label}| command and cross references
to them are made with the \verb|\ref{label}| command.
This an example cross-reference to Equation~\ref{eq:example}. You can also write equations inline, like this: $A=\pi r^2$.


% \section*{Limitations}

\section*{Team Contributions}

Write in this section a few sentences describing the contributions of each team member. What did each member work on? Refer to item 7.

% Bibliography entries for the entire Anthology, followed by custom entries
%\bibliography{custom,anthology-overleaf-1,anthology-overleaf-2}

% Custom bibliography entries only
\bibliography{custom}

% \appendix

% \section{Example Appendix}
% \label{sec:appendix}

% This is an appendix.

\end{document}
